\SourceLeaf{435}%
\[
\begin{aligned}
={}&\int\frac{\dfrac{dx}{dw}.d\delta x+\dfrac{dy}{dw}.d\delta y+\dfrac{dz}{dw}.d\delta z}{\BareRoot{\left\{\left(\dfrac{dx}{dw}\right)^{2}+\left(\dfrac{dy}{dw}\right)^{2}+\left(\dfrac{dz}{dw}\right)^{2}\right\}}}\\[5pt]
={}&\frac{\dfrac{dx}{dw}.\delta x+\dfrac{dy}{dw}.\delta y+\dfrac{dz}{dw}.\delta z}{\BareRoot{\left\{\left(\dfrac{dx}{dw}\right)^{2}+\left(\dfrac{dy}{dw}\right)^{2}+\left(\dfrac{dz}{dw}\right)^{2}\right\}}}\\[5pt]
&-\int\Biggl\{\delta x.d\frac{\dfrac{dx}{dw}}{\BareRoot{\left\{\left(\dfrac{dx}{dw}\right)^{2}+\left(\dfrac{dy}{dw}\right)^{2}+\left(\dfrac{dz}{dw}\right)^{2}\right\}}}\\[5pt]
&\qquad+\delta y.d\frac{\dfrac{dy}{dw}}{\BareRoot{\left\{\left(\dfrac{dx}{dw}\right)^{2}+\left(\dfrac{dy}{dw}\right)^{2}+\left(\dfrac{dz}{dw}\right)^{2}\right\}}}\\[5pt]
&\qquad+\delta z.d\frac{\dfrac{dz}{dw}}{\BareRoot{\left\{\left(\dfrac{dx}{dw}\right)^{2}+\left(\dfrac{dy}{dw}\right)^{2}+\left(\dfrac{dz}{dw}\right)^{2}\right\}}}\Biggr\}.
\end{aligned}
\]

Was hier unter dem Integrationszeichen steht, muss bekanntlich für den Fall des Minimum verschwinden; in so fern nun die krumme Linie sich auf einer gegebenen krummen Fläche befinden soll, deren Gleichung
\[
Pdx+Qdy+Rdz=0
\]
ist, muss auch zwischen den Variationen $\delta x$, $\delta y$, $\delta z$ die Gleichung
\[
P\delta x+Q\delta y+R\delta z=0
\]
statthaben, woraus man nach bekannten Principien leicht schliesst, dass die Differentiale
\[
\begin{gathered}
d.\frac{\dfrac{dx}{dw}}{\BareRoot{\left\{\left(\dfrac{dx}{dw}\right)^{2}+\left(\dfrac{dy}{dw}\right)^{2}+\left(\dfrac{dz}{dw}\right)^{2}\right\}}},\\[5pt]
d.\frac{\dfrac{dy}{dw}}{\BareRoot{\left\{\left(\dfrac{dx}{dw}\right)^{2}+\left(\dfrac{dy}{dw}\right)^{2}+\left(\dfrac{dz}{dw}\right)^{2}\right\}}},\\[5pt]
d.\frac{\dfrac{dz}{dw}}{\BareRoot{\left\{\left(\dfrac{dx}{dw}\right)^{2}+\left(\dfrac{dy}{dw}\right)^{2}+\left(\dfrac{dz}{dw}\right)^{2}\right\}}}
\end{gathered}
\]
resp. den Grössen $P,Q,R$ proportional sein müssen. Ist $ds$ ein Element der Curve, $\lambda$ der Punkt auf der Hülfskugel, der die Richtung dieses Elements, $L$ der Punkt, der wie oben die Richtung der Normale ausdrückt, und sind $\xi,\eta,\zeta$; $X,Y,Z$ die Coordinaten der Punkte $\lambda,L$ in Beziehung auf den Mittelpunkt der Hülfskugel, dann hat man:
\[
\begin{gathered}
dx=\xi ds,\qquad dy=\eta ds,\qquad dz=\zeta ds,\\
\xi\xi+\eta\eta+\zeta\zeta=1;
\end{gathered}
\]
man sieht also, dass die obigen Differentiale $=d\xi$, $d\eta$, $d\zeta$ werden. Und da $P,Q,R$ den Grössen $X,Y,Z$ proportional sind, so ist der Character der kürzesten Linie der, dass
\[
\frac{d\xi}{X}=\frac{d\eta}{Y}=\frac{d\zeta}{Z}.
\]
